% Chapter 8 Introduction (untitled, ~1 page)
% Bridges from Parts II–III to Part IV: cross-correlations at cosmological scales

The preceding chapters looked for localized dark matter signatures: the concentrated gamma-ray emission expected from the Galactic Center --- where the dark matter density reaches its peak --- and from individual subhalos that remain unresolved among the unassociated sources detected by the \textit{Fermi} Large Area Telescope.
Neither strategy returned a definitive signal.
The Galactic Center excess persists, but a decade of increasingly sophisticated analyses has not resolved the fundamental degeneracy between a dark matter annihilation origin and a population of unresolved millisecond pulsars~\cite{Daylan:2014rsa,Bartels:2015aea,Leane:2019xiy,Chang:2019ars}. The search for dark matter subhalos among unassociated sources~\cite{Amerio:2025fhz} found no significant evidence for an anomalous population at current sensitivity levels, consistent with theoretical expectations --- a result that establishes competitive constraints but does not close the door. 
We then moved to studying the unresolved gamma-ray sky. 
At the population level, the source-count distribution and probabilistic catalog constructed in Chapters~\ref{ch:6} and \ref{ch:7} mapped the unresolved gamma-ray sky to fluxes roughly 50 times below the detection threshold~\cite{Amerio:2023uet,Amerio:2023rjn,Cuoco-1pdf}, characterizing the astrophysical landscape in which any dark matter contribution must reside, but without working on isolating a signal.

% This chapter opens Part~IV by stepping back from individual sources and asking a different question: if dark matter annihilates or decays, it does so inside every gravitationally bound halo and subhalo across all cosmic epochs.
In this chapter we focus on how to isolate a signal from dark matter hidden in the \UGRB.
In fact, if dark matter annihilates or decays, it does so inside every gravitationally bound structure across all cosmic epochs,
and the accumulated, unresolved emission from this entire population would contribute to the \UGRB.
Separating a dark matter component from the astrophysical foreground --- blazars, star-forming galaxies, misaligned active galactic nuclei --- requires more than spectral information alone, since distinct source classes partially overlap in energy.
%
What discriminates dark matter is its unique spatial distribution: the large-scale structure of the universe, traced by dark matter halos from galaxy clusters down to sub-Earth-mass clumps, imprints specific angular anisotropy patterns on the gamma-ray sky that differ systematically from those of astrophysical backgrounds~\cite{DGRB-review,Fornengo:2013rga,Camera:2012cj}.
%
The angular cross-power spectrum between gamma-ray maps and a gravitational tracer --- such as a galaxy catalog or a cosmic shear field --- provides a statistically powerful handle on this signal, because it vanishes for astrophysical source populations at redshifts incompatible with the tracer, but is maximal for dark matter emission that clusters with local structure.

We now examine this cross-correlation approach in detail from three directions.
Section~\ref{sec:8.1} places the technique in the context of the thesis as a whole, tracing the conceptual arc from resolved Galactic targets to the cosmic web and establishing why a shift in scale is both necessary and well-motivated given the null results of the previous searches. Section~\ref{sec:8.2} develops the physical and mathematical principles of the cross-correlation method --- the role of window functions and tomographic redshift selectivity, the noise cancellation advantage over auto-correlations, and the identification of low-redshift galaxy catalogs as the optimal gravitational tracer for dark matter annihilation. Section~\ref{sec:8.3} introduces the Cherenkov Telescope Array Observatory (CTAO), which extends the accessible energy range an order of magnitude above the \textit{Fermi}-LAT ceiling and brings effective collecting areas six orders of magnitude larger, opening the TeV mass window to cross-correlation searches for the first time. The paper that follows then presents a full sensitivity forecast for detecting dark matter signals through cross-correlations between the CTAO extragalactic survey and the 2MASS galaxy catalog~\cite{Pinetti:2025hgd}.
